\documentclass[a4paper]{article}

\usepackage[english]{babel}
%\usepackage[utf8x]{inputenc}
\usepackage{amsmath}
\usepackage{graphicx}
\usepackage[colorinlistoftodos]{todonotes}
\usepackage[all]{xy}
\usepackage{dsfont}
\usepackage{amssymb}
\usepackage{bbold}
\usepackage{amsthm}

\usepackage{tikz}
\usepackage{verbatim}
\usetikzlibrary{matrix}
\usetikzlibrary{cd}

\usepackage[autostyle=true]{csquotes}
\usepackage{pdfpages}
\usepackage{hyperref}
\usepackage{fancyhdr}



\newtheorem{theorem}{Theorem}[section]
\newtheorem{proposition}{Proposition}[theorem]
\newtheorem{lemma}{Lemma}[theorem]
\newtheorem{corollary}{Corollary}[theorem]
\newtheorem*{definition}{Definition}


\newcommand{\norm}[1]{\left\lVert#1\right\rVert}

\newcommand{\Rn}{\mathbb{R}^n}


\newcommand{\too}[1]{\overset{#1}{\longrightarrow}}
\newcommand{\hookar}[1]{\overset{#1}{\hookrightarrow}}
\newcommand{\colim}[1]{ \underset{#1}{\texttt{colim}} \text{ } }
\newcommand{\fset}[1]{ \{ #1_0 , \dots , #1_n \} }
\newcommand{\ddt}{\frac{d}{dt}}


\newcommand{\N}{\mathbb{N}}
\newcommand{\Z}{\mathbb{Z}}







\fancyhead[L]{\includegraphics[scale=0.1]{UFT.png}}

\begin{document}


\thispagestyle{fancy}
\vspace*{1cm}

\begin{minipage}[r]{12cm}
    \begin{flushright}
\begin{tikzpicture}
\draw[green, very thick] (0,0) rectangle (3,2);
\end{tikzpicture}
    \end{flushright}
\end{minipage}

\begin{center}
\Large{\textbf{Control Recuperativo\\ Álgebra Lineal\\ }}
\end{center}
\vspace{1cm}
\begin{minipage}[r]{15cm}
\begin{flushleft}

Fecha: 23 de junio 2025 \\%
Duración: 120 minutos  \\

\end{flushleft}
\end{minipage}
\vspace{2cm}
\begin{center}
\Large{\textbf{Facultad de Ingeniería}}
\end{center}
\vspace{1cm}
\textbf{En caso de que un alumno sea sorprendido en actitudes deshonestas, ya sea intentando o
dejándose copiar, será calificado con nota mínima 1.0 y perderá toda posibilidad de eliminar
dicha nota. Además, esta conducta quedará archivada como antecedente en su carpeta y
podrá ser causal de eliminación.
Así mismo, está prohibido salir de la sala, usar o mantener sobre la mesa teléfonos celulares
o cualquier tipo de aparato electrónico, a no ser que el profesor lo requiera.}\\



\vspace{5cm}
\makebox[12cm]{\textbf{Nombre Alumno}:\ \hrulefill} 
\newpage
\thispagestyle{empty}
\section*{Preguntas}
\begin{enumerate}



\item {\bf (15 puntos)} Dado el siguiente sistema de ecuaciones lineales:

    \begin{align*}
 x + 3y &= 1 \\
 x + 2y +z &= 0\\
2 x + 5y +(\alpha -1) z &= 1-\alpha 
\end{align*}

    tenga solución única, infinitas soluciones o ninguna solución. Para $\alpha=0$, determine el conjunto soluci\'on del sistema.



 \textbf{Soluci\'on:}\\
Escribimos la matriz aumentada del sistema:

\[
\left[
\begin{array}{ccc|c}
1 & 3 & 0 & 1 \\
1 & 2 & 1 & 0 \\
2 & 5 & \alpha - 1 & 1 - \alpha
\end{array}
\right]
\]

Restamos la fila 1 a la fila 2:

\[
F_2 \leftarrow F_2 - F_1:
\left[
\begin{array}{ccc|c}
1 & 3 & 0 & 1 \\
0 & -1 & 1 & -1 \\
2 & 5 & \alpha - 1 & 1 - \alpha
\end{array}
\right]
\]

Luego, restamos 2 veces la fila 1 a la fila 3:

\[
F_3 \leftarrow F_3 - 2F_1:
\left[
\begin{array}{ccc|c}
1 & 3 & 0 & 1 \\
0 & -1 & 1 & -1 \\
0 & -1 & \alpha - 1 & -1 - \alpha
\end{array}
\right]
\]

Restamos la fila 2 a la fila 3:

\[
F_3 \leftarrow F_3 - F_2:
\left[
\begin{array}{ccc|c}
1 & 3 & 0 & 1 \\
0 & -1 & 1 & -1 \\
0 & 0 & \alpha - 2 & -\alpha
\end{array}
\right]
\]


\begin{itemize}
\item Si \( \alpha - 2 \neq 0 \), es decir, \( \alpha \neq 2 \), entonces hay una única solución porque el sistema es compatible determinado (rango de la matriz de coeficientes y el de la aumentada es 3).
\item Si \( \alpha = 2 \), la ecuación implicar\'ia:
\[
0 = -2
\]
lo cual es una contradicción. Por lo tanto, el sistema es incompatible: \textbf{no tiene solución}.
\end{itemize}



Sustituimos \( \alpha = 0 \) en el sistema original:
\[
\left[
\begin{array}{ccc|c}
1 & 3 & 0 & 1 \\
0 & -1 & 1 & -1 \\
0 & 0 & - 2 & 0
\end{array}
\right]
\]


\textbf{Por lo tanto, la única solución es:}
\[
\boxed{(x, y, z) = (-2,\ 1,\ 0)}
\]



\newpage
    \item Considere el conjunto $S=\{ ( 0 , 1 , 2 ), ( 1 , 1 , 2 ), ( 0 , 0 , -1) \}$.
    \begin{enumerate}
        \item {\bf (3 puntos)} Muestre que $S$ es base de $\mathbb{R}^{3}$.
        \item {\bf (9 puntos)} Encuentre una base ortonormal de $\mathbb{R}^{3}$ a partir de esta base usando el procedimiento de Gram-Schmidt.
        \item {\bf (3 puntos)} Encuentre las coordenadas de $\{( -1 ,0 , 1 ) \}$ en esta base.    

\end{enumerate} 
\textbf{Solución:}
\begin{enumerate}

\item Muestre que \( S \) es base de \( \mathbb{R}^3 \).

Para probar que \( S \) es una base de \( \mathbb{R}^3 \), basta demostrar que los vectores son linealmente independientes.

Formamos la matriz con los vectores como columnas:

\[
M = \begin{pmatrix}
0 & 1 & 0 \\
1 & 1 & 0 \\
2 & 2 & -1
\end{pmatrix}
\]

Calculamos su determinante:

\[
\det(M) = 
\begin{vmatrix}
0 & 1 & 0 \\
1 & 1 & 0 \\
2 & 2 & -1
\end{vmatrix}
= 0 \cdot \begin{vmatrix} 1 & 0 \\ 2 & -1 \end{vmatrix}
- 1 \cdot \begin{vmatrix} 1 & 0 \\ 2 & -1 \end{vmatrix}
+ 0 \cdot \begin{vmatrix} 1 & 1 \\ 2 & 2 \end{vmatrix}
= - (1 \cdot (-1) - 0 \cdot 2) = -(-1) = 1
\]

Como el determinante es distinto de cero, los vectores son linealmente independientes, por lo tanto, al ser 3 vectores en espacio de dimensión 3, éstos forman una base de \( \mathbb{R}^3 \).

\vspace{0.4cm}

\item Encuentre una base ortonormal de \( \mathbb{R}^3 \) a partir de esta base usando el procedimiento de Gram-Schmidt.

Sea
\[
\vec{v}_1 = (0,1,2), \quad \vec{v}_2 = (1,1,2), \quad \vec{v}_3 = (0,0,-1)
\]

Definimos \( \vec{w}_1 = \vec{v}_1 \)
\[
\vec{w}_1 = (0,1,2), \quad
\|\vec{w}_1\| = \sqrt{1^2 + 2^2} = \sqrt{5}
\Rightarrow \vec{u}_1 = \frac{1}{\sqrt{5}}(0,1,2)
\]


Ortogonalizamos \( \vec{v}_2 \):

\[
\text{proy}_{\vec{w}_1}(\vec{v}_2) = \frac{\langle \vec{v}_2, \vec{w}_1 \rangle}{\langle \vec{w}_1, \vec{w}_1 \rangle} \vec{w}_1
= \frac{(1)(0)+(1)(1)+(2)(2)}{1+4} (0,1,2) = \frac{5}{5}(0,1,2) = (0,1,2)
\]

\[
\vec{w}_2 = \vec{v}_2 - \text{proy}_{\vec{w}_1}(\vec{v}_2) = (1,1,2) - (0,1,2) = (1,0,0)
\Rightarrow \|\vec{w}_2\| = 1, \quad \vec{u}_2 = (1,0,0)
\]

Ortogonalizamos \( \vec{v}_3 \):

\[
\text{proy}_{\vec{u}_1}(\vec{v}_3) = \frac{\langle \vec{v}_3, \vec{u}_1 \rangle}{\langle \vec{u}_1, \vec{u}_1 \rangle} \vec{u}_1
= \frac{(0)(0)+(0)(1)+(-1)(2)}{5}(0,1,2) = \frac{-2}{5}(0,1,2) = (0,-\frac{2}{5},-\frac{4}{5})
\]

\[
\text{proy}_{\vec{u}_2}(\vec{v}_3) = \langle \vec{v}_3, \vec{u}_2 \rangle \vec{u}_2 = (0)(1)+(0)(0)+(-1)(0) = 0
\Rightarrow \text{proyección total} = (0,-\frac{2}{5},-\frac{4}{5})
\]

\[
\vec{u}_3 = \vec{v}_3 - \text{proy} = (0,0,-1) - (0,-\frac{2}{5},-\frac{4}{5}) = (0,\frac{2}{5},-\frac{1}{5})
\]

\[
\|\vec{u}_3\| = \sqrt{\left(\frac{2}{5}\right)^2 + \left(-\frac{1}{5}\right)^2} = \sqrt{\frac{4}{25} + \frac{1}{25}} = \sqrt{\frac{5}{25}} = \frac{1}{\sqrt{5}}
\Rightarrow \vec{e}_3 = \sqrt{5} \cdot (0,\frac{2}{5},-\frac{1}{5}) = (0,\frac{2}{\sqrt{5}}, -\frac{1}{\sqrt{5}})
\]

\textbf{Resultado:} La base ortonormal es:
\[
\left\{
\frac{1}{\sqrt{5}}(0,1,2),\;
(1,0,0),\;
\frac{1}{\sqrt{5}}(0,2,-1)\right)
\right\}
\]

\vspace{0.4cm}

\item Encuentre las coordenadas de \( \vec{x} = (-1, 0, 1) \) en esta base.

Sea \( B = \{\vec{v}_1, \vec{v}_2, \vec{v}_3\} \), queremos encontrar escalares \( a, b, c \) tales que:

\[
(-1,0,1) = a(0,1,2) + b(1,1,2) + c(0,0,-1)
\]

Para encontrar las coordenadas \( (a, b, c) \) de \( \vec{x} = (-1,0,1) \) en esta base ortonormal, usamos:

\[
a = \langle \vec{x}, \vec{u}_1 \rangle, \quad
b = \langle \vec{x}, \vec{u}_2 \rangle, \quad
c = \langle \vec{x}, \vec{u}_3 \rangle
\]

\textbf{Cálculo:}

\[
a = \left\langle (-1, 0, 1), \frac{1}{\sqrt{5}}(0,1,2) \right\rangle
= \frac{1}{\sqrt{5}} \left( 0\cdot(-1) + 1\cdot0 + 2\cdot1 \right) = \frac{2}{\sqrt{5}}
\]

\[
b = \langle (-1, 0, 1), (1,0,0) \rangle = (-1)\cdot1 + 0\cdot0 + 1\cdot0 = -1
\]

\[
c = \left\langle (-1, 0, 1), \left(0, \frac{2}{\sqrt{5}}, -\frac{1}{\sqrt{5}}\right) \right\rangle
= 0\cdot(-1) + 0\cdot\frac{2}{\sqrt{5}} + 1\cdot\left(-\frac{1}{\sqrt{5}}\right) = -\frac{1}{\sqrt{5}}
\]

\textbf{Respuesta:}  
Las coordenadas de \( \vec{x} = (-1,0,1) \) en la base ortonormal son:

\[
\left( \frac{2}{\sqrt{5}},\; -1,\; -\frac{1}{\sqrt{5}} \right)
\]
\end{enumerate}

\newpage
    
    \item Sea $H$ el subconjunto de $\mathbb{R}^3$ definido por
\[H=\{(x,y,z)\in\mathbb{R}^3\colon 2x-y=0  \}\]

\begin{enumerate}
\item {\bf (3 puntos)} Determine si los vectores $( 1 ,2 , 3 ) $ y $ ( -1 ,0 , -2 ) $ pertenecen a $H$.
   \item {\bf (6 puntos)} Demuestre que $H$ es un subespacio vectorial de $\mathbb{R}^3$.
    \item {\bf (6 puntos)} Encuentre una base de $H$ y determine su dimensión.

\end{enumerate}
\newpage
\phantom{0}
\newpage
    \item Considere la siguiente matriz
\[    A=\begin{pmatrix} -1 &1  \\ 0 & 1/2 \end{pmatrix}.\]
\begin{itemize}
\item[(a)] {\bf (10 puntos)} Obtenga los valores propios de \(A\) y sus vectores propios asociados.
\item[(b)] {\bf (5 puntos)} Presente, si es posible, la diagonalización de esta matriz, encontrando matrices \(P\) y \(D\) que cumplan
\[
A = P D P^{-1}
\]
\end{itemize}
\end{enumerate}
\textbf{Solución:}

\begin{itemize}

\item[(a)] \textbf{(10 puntos)} Obtenga los valores propios de \(A\) y sus vectores propios asociados.

Para hallar los valores propios, resolvemos el polinomio característico:

\[
\det(A - \lambda I) = \det\left( \begin{pmatrix} -1 - \lambda & 1 \\ 0 & \frac{1}{2} - \lambda \end{pmatrix} \right)
= (-1 - \lambda)\left( \frac{1}{2} - \lambda \right)
\]

\[
\Rightarrow (-1 - \lambda)\left( \frac{1}{2} - \lambda \right) = 0
\]

\[
\Rightarrow \lambda_1 = -1, \quad \lambda_2 = \frac{1}{2}
\]

\vspace{0.3cm}

\textbf{Para } \( \lambda_1 = -1 \):

Resolvemos \( (A + I)\vec{v} = 0 \):

\[
A + I = \begin{pmatrix} 0 & 1 \\ 0 & \frac{3}{2} \end{pmatrix}
\Rightarrow \begin{pmatrix} 0 & 1 \\ 0 & \frac{3}{2} \end{pmatrix} \begin{pmatrix} x \\ y \end{pmatrix} = \begin{pmatrix} 0 \\ 0 \end{pmatrix}
\]

\[
\Rightarrow y = 0,\quad x\text{ libre}
\Rightarrow \vec{v}_1 = \begin{pmatrix} 1 \\ 0 \end{pmatrix}
\]

\vspace{0.3cm}

\textbf{Para } \( \lambda_2 = \frac{1}{2} \):

\[
A - \frac{1}{2}I = \begin{pmatrix} -\frac{3}{2} & 1 \\ 0 & 0 \end{pmatrix}
\Rightarrow \begin{pmatrix} -\frac{3}{2} & 1 \\ 0 & 0 \end{pmatrix} \begin{pmatrix} x \\ y \end{pmatrix} = \begin{pmatrix} 0 \\ 0 \end{pmatrix}
\]

\[
\Rightarrow -\frac{3}{2}x + y = 0 \Rightarrow y = \frac{3}{2}x
\Rightarrow \vec{v}_2 = \begin{pmatrix} 2 \\ 3 \end{pmatrix}
\]

\vspace{0.5cm}

\end{itemize}
\newpage


\section*{Formulario}
\begin{enumerate}
    \item \textbf{Coordenada de $v$ respecto de la base ortogonal $\{w_i\}$:} $a_i=\frac{\langle v, w_i \rangle}{\langle w_i, w_i \rangle}$.


    \item \textbf{Vector normal: } vector de norma 1. Vector normal asociado a $v\in V$ es $\hat{v}=\dfrac{v}{\|v\|}.$



    \item \textbf{Proceso de ortogonalizaci\'on de Gram Schmidt.}

Sea $B = \{v_i\}_{i=1}^m $ una base ordenada cualquiera de $W $.
\begin{enumerate}
\item Escriba $w_1 = v_1$. Calcule $w_2 = v_2 - \dfrac{\langle v_2,v_1\rangle}{\langle v_1,v_1\rangle}v_1$.
%\item Suponemos que esta construcción de $w_k$ es posible hacerla hasta
%obtener el conjunto ortogonal $ \{w_i\}_{i=1}^{m-1}$ con $w_k=v_k-\sum_{i=1}^{k-1}c_iw_i$.
\item Calcule de manera sucesiva
\[w_i=v_i-\dfrac{\langle v_i,w_1\rangle}{\langle w_1,w_1\rangle}w_1-\dfrac{\langle v_i,w_2\rangle}{\langle w_2,w_2\rangle}w_2-\dots-\dfrac{\langle v_i,w_{i-1}\rangle}{\langle w_{i-1},w_{i-1}\rangle}w_{i-1}.\]
Obteniendo un conjunto ortogonal.
\item Para conseguir un conjunto ortonormal, calcule $u_i=\dfrac{w_i}{\|w_i\|}$.
\end{enumerate}

\end{enumerate}



\end{document}